% Options for packages loaded elsewhere
\PassOptionsToPackage{unicode}{hyperref}
\PassOptionsToPackage{hyphens}{url}
%
\documentclass[
]{book}
\usepackage{amsmath,amssymb}
\usepackage{lmodern}
\usepackage{iftex}
\ifPDFTeX
  \usepackage[T1]{fontenc}
  \usepackage[utf8]{inputenc}
  \usepackage{textcomp} % provide euro and other symbols
\else % if luatex or xetex
  \usepackage{unicode-math}
  \defaultfontfeatures{Scale=MatchLowercase}
  \defaultfontfeatures[\rmfamily]{Ligatures=TeX,Scale=1}
\fi
% Use upquote if available, for straight quotes in verbatim environments
\IfFileExists{upquote.sty}{\usepackage{upquote}}{}
\IfFileExists{microtype.sty}{% use microtype if available
  \usepackage[]{microtype}
  \UseMicrotypeSet[protrusion]{basicmath} % disable protrusion for tt fonts
}{}
\makeatletter
\@ifundefined{KOMAClassName}{% if non-KOMA class
  \IfFileExists{parskip.sty}{%
    \usepackage{parskip}
  }{% else
    \setlength{\parindent}{0pt}
    \setlength{\parskip}{6pt plus 2pt minus 1pt}}
}{% if KOMA class
  \KOMAoptions{parskip=half}}
\makeatother
\usepackage{xcolor}
\usepackage{longtable,booktabs,array}
\usepackage{calc} % for calculating minipage widths
% Correct order of tables after \paragraph or \subparagraph
\usepackage{etoolbox}
\makeatletter
\patchcmd\longtable{\par}{\if@noskipsec\mbox{}\fi\par}{}{}
\makeatother
% Allow footnotes in longtable head/foot
\IfFileExists{footnotehyper.sty}{\usepackage{footnotehyper}}{\usepackage{footnote}}
\makesavenoteenv{longtable}
\usepackage{graphicx}
\makeatletter
\def\maxwidth{\ifdim\Gin@nat@width>\linewidth\linewidth\else\Gin@nat@width\fi}
\def\maxheight{\ifdim\Gin@nat@height>\textheight\textheight\else\Gin@nat@height\fi}
\makeatother
% Scale images if necessary, so that they will not overflow the page
% margins by default, and it is still possible to overwrite the defaults
% using explicit options in \includegraphics[width, height, ...]{}
\setkeys{Gin}{width=\maxwidth,height=\maxheight,keepaspectratio}
% Set default figure placement to htbp
\makeatletter
\def\fps@figure{htbp}
\makeatother
\setlength{\emergencystretch}{3em} % prevent overfull lines
\providecommand{\tightlist}{%
  \setlength{\itemsep}{0pt}\setlength{\parskip}{0pt}}
\setcounter{secnumdepth}{5}
\usepackage{booktabs}
\ifLuaTeX
  \usepackage{selnolig}  % disable illegal ligatures
\fi
\usepackage[]{natbib}
\bibliographystyle{plainnat}
\IfFileExists{bookmark.sty}{\usepackage{bookmark}}{\usepackage{hyperref}}
\IfFileExists{xurl.sty}{\usepackage{xurl}}{} % add URL line breaks if available
\urlstyle{same} % disable monospaced font for URLs
\hypersetup{
  pdftitle={통계와 분석},
  pdfauthor={류춘렬},
  hidelinks,
  pdfcreator={LaTeX via pandoc}}

\title{통계와 분석}
\author{류춘렬}
\date{2023-02-17}

\begin{document}
\maketitle

{
\setcounter{tocdepth}{1}
\tableofcontents
}
\hypertarget{uxc548uxb0b4}{%
\chapter*{안내}\label{uxc548uxb0b4}}
\addcontentsline{toc}{chapter}{안내}

이 책은 \textbf{통계와 분석}의 강의노트입니다. 이 과목 수강 이외에 이 책을 사용하는 목적은 저자와 상의해야 합니다.

류춘렬, cryu@kookmin.ac.kr

이 강의에서 배우는 것은 아래와 같습니다. R을 적극적으로 활용한다는 점에서 데이터 사이언스 분야와 접근이 같지만 사회과학적 통계 이론과 데이터 분석에도 중점을 두므로 두 분야의 균형을 이룹니다.

\begin{itemize}
\item
  데이터 다루기: 가져오기, 내보내기, 재구성, 변환
\item
  탐색적 데이터 분석: 데이터 시각화, 다변인 통계 기법
\item
  통계적 추론: 가설을 검정하고 신뢰 구간을 계산하는 현대적 기법
\item
  예측적 모델 수립: 예측과 분류, 전망을 위한 회귀 모델과 기계 학습 기법
\item
  시뮬레이션: 표본 크기 계산과 통계 기법의 평가를 위한 방법
\end{itemize}

데이터 사이언스 분야는 Thulin(2021)을 많이 참고했습니다.

\url{https://modernstatisticswithr.com/introduction.html}

이 노트는 입말로 써 있습니다. 내가 생각하기에, 입말이 글말보다 이해가 빠릅니다. 처음은 어색하지만 익숙해지면 편합니다.

한 학기 동안 잘 달려갑시다!

참고로

이 책은 \textbf{RMarkdown}과 \textbf{Bookdown}을 이용해서 작성하였습니다.

\hypertarget{uxac15uxc758-uxc18cuxac1c}{%
\chapter{강의 소개}\label{uxac15uxc758-uxc18cuxac1c}}

\hypertarget{ruxc758-uxc138uxc0c1uxc73cuxb85c}{%
\section{R의 세상으로}\label{ruxc758-uxc138uxc0c1uxc73cuxb85c}}

R은 여러 통계 프로그램과 다릅니다. 공짜고 효율적이며 빠르고 최신이에요. 거대한 이용자 공동체가 서로 도우며 R을 개선하고 있습니다. R은 그 어떤 통계 프로그램보다 더 많은 분석 기법을 포함하고 있어서 다양한 학문 분야에서 사용합니다. 소통학이나 정치학, 천문학, 의학 등등. 유사한 분석을 반복하기도 아주 쉬워요.

R은 통계학자가 개발해서 데이터를 다루기가 여타 프로그램보다 편합니다. 여러 면에서 현존 최강이라고 불러도 될 정도지요.

\hypertarget{uxc0acuxc804-uxc9c0uxc2dd}{%
\section{사전 지식}\label{uxc0acuxc804-uxc9c0uxc2dd}}

이 강의를 듣는데 통계에 관한 사전 지식은 필요 없어요. 원래 초보자를 위한 것이라 차근차근하게 밑바닥부터 나갈 거에요. 물론 컴퓨터 프로그래밍에 대한 사전 지식은 더 필요 없지요.

처음부터 끝까지 차례대로 공부할 필요는 없어요. 중간에 건너 뛰어려면 그래도 됩니다. 그런데 R을 배우려면 R을 직접 사용해보는 게 필요하니까 해당 세션의 R 실습 문제는 꼭 풀어보는 게 좋아요.

\hypertarget{uxade0uxd615}{%
\section{균형}\label{uxade0uxd615}}

R을 이용해서 통계에 접근하는 길은 대체로 두 가지로 나누어져요. 하나는 사회과학의 전통적 통계를 공부하는 것이고 다른 하나는 요즘 뜨는 데이터 사이언스를 공부하는 거에요. 이 강의에서는 양쪽을 다 다루지만 그래도 사회과학의 전통적 통계 쪽으로 기울기는 할 겁니다. 이게 Thulin과 이 강의가 차이가 생기는 이유입니다.

  \bibliography{book.bib,packages.bib}

\end{document}
